\documentclass[11pt,a4paper]{article}
\usepackage[margin=2.5cm]{geometry}
\usepackage{graphicx,hyperref,xcolor,listings,booktabs,fancyhdr,titlesec}
\usepackage[T1]{fontenc}
\usepackage{lmodern}

\definecolor{cyber}{HTML}{00FFD6}
\definecolor{bg}{HTML}{0A0A0A}
\definecolor{mg}{HTML}{FF66AC}
\hypersetup{colorlinks,linkcolor=cyber,urlcolor=cyber}

\pagestyle{fancy}
\fancyhf{}
\fancyhead[L]{\textcolor{cyber}{\textbf{FLLC CYBERDECK}}}
\fancyhead[R]{\textcolor{mg}{\thepage}}
\fancyfoot[C]{\small\textcolor{gray}{Authorized Hardware Analysis Platform --- Internal Use Only}}
\renewcommand{\headrulewidth}{0.4pt}

\titleformat{\section}{\Large\bfseries\color{cyber}}{\thesection}{1em}{}
\titleformat{\subsection}{\large\bfseries\color{mg}}{\thesubsection}{1em}{}

\lstset{basicstyle=\ttfamily\small, backgroundcolor=\color{black!5},
        frame=single, rulecolor=\color{gray}, breaklines=true}

\title{%
  \vspace{-1cm}
  {\Huge\bfseries\textcolor{cyber}{FLLC CyberDeck}}\\[6pt]
  {\Large Architecture \& Design Document}\\[4pt]
  {\normalsize Version 2.0 --- \today}
}
\author{FLLC Engineering Division}
\date{}

\begin{document}
\maketitle
\thispagestyle{fancy}
\tableofcontents
\newpage

%----------------------------------------------------------------------
\section{Executive Overview}
FLLC CyberDeck is a local-first, Docker-composed hardware-analysis workbench
designed for the Raspberry Pi~400 CyberDeck build (DFRobot 3.5'' TFT,
TP-Link TL-WN722N, Adafruit enclosure), small x86 workstations, and internal bench
appliances.  It provides:
\begin{itemize}
  \item Project / target / session / evidence tracking
  \item Hardware interface documentation (UART, I2C, SPI, JTAG/SWD, eMMC)
  \item Passive IoT inventory with CSV import/export
  \item JWT-authenticated multi-user access
  \item Local AI assistance via Ollama or OpenAI-compatible endpoints
  \item Evidence artifact upload with SHA-256 integrity
  \item Printable HTML reports with severity matrix and browser PDF export
  \item GitHub Actions CI/CD with container publishing to GHCR
\end{itemize}

%----------------------------------------------------------------------
\section{System Architecture}

\subsection{Component Diagram}
\begin{verbatim}
+-----------+      +------------+      +------------+
|  Web UI   |----->|  FastAPI   |----->|   Worker   |
| Express   | :3000|  API :8000 |      | Python bg  |
+-----------+      +-----+------+      +------+-----+
                         |                    |
                   +-----v----+         +-----v-----+
                   | SQLite / |         |  /data/    |
                   | state.json|        |  reports/  |
                   +----------+         |  uploads/  |
                                        +-----------+
                   +--------------+
                   | Ollama :11434|  (optional)
                   +--------------+
\end{verbatim}

\subsection{Service Inventory}
\begin{tabular}{llll}
\toprule
\textbf{Service} & \textbf{Tech} & \textbf{Port} & \textbf{Role}\\
\midrule
\texttt{api}    & FastAPI / Python 3.11 & 8000 & REST API, auth, state \\
\texttt{web}    & Express / Node 18     & 3000 & UI proxy, static HTML  \\
\texttt{worker} & Python 3.11           & ---  & Background report gen  \\
\texttt{ollama} & Ollama (optional)     & 11434& Local LLM inference    \\
\bottomrule
\end{tabular}

%----------------------------------------------------------------------
\section{Authentication}
\begin{itemize}
  \item \textbf{Mechanism}: JWT bearer tokens via \texttt{/auth/login}
  \item \textbf{Default credentials}: \texttt{admin / cyberdeck} (change immediately)
  \item \textbf{Token lifetime}: 24 hours
  \item \textbf{Storage}: Credentials hashed with SHA-256 in state file
  \item \textbf{Middleware}: \texttt{verify\_token()} dependency on protected routes
\end{itemize}

%----------------------------------------------------------------------
\section{Data Model}
\begin{tabular}{ll}
\toprule
\textbf{Entity} & \textbf{Key Fields} \\
\midrule
Project         & id, name, summary, classification, authorization\_scope \\
DeviceTarget    & id, project\_id, name, kind, vendor, model, firmware\_version \\
HardwareSession & id, project\_id, target\_id, interface\_type, voltage, pin\_map \\
EvidenceArtifact& id, project\_id, kind, path, sha256, notes \\
Report          & id, project\_id, title, state, findings[], generated\_at \\
AgentTask       & id, project\_id, mode, state \\
\bottomrule
\end{tabular}

%----------------------------------------------------------------------
\section{AI Integration}
Three adapters are supported:
\begin{enumerate}
  \item \textbf{NoneAdapter} --- passthrough, app works fully
  \item \textbf{OllamaAdapter} --- real HTTP calls to \texttt{/api/generate},
        health-checked at \texttt{/api/tags}
  \item \textbf{OpenAICompatibleAdapter} --- any \texttt{/v1/chat/completions}
        endpoint
\end{enumerate}
AI is \textbf{never} used for payload generation or exploitation.

%----------------------------------------------------------------------
\section{Reporting}
Reports are generated as self-contained HTML with:
\begin{itemize}
  \item Title page with project metadata
  \item Authorization statement
  \item Executive summary (optionally AI-assisted)
  \item Findings table with severity matrix (CRITICAL / HIGH / MEDIUM / LOW / INFO)
  \item Evidence references with SHA-256 hashes
  \item Analyst notes, recommendations, appendix
  \item Print CSS and \texttt{window.print()} button for browser PDF export
\end{itemize}

%----------------------------------------------------------------------
\section{File Upload}
\begin{itemize}
  \item Endpoint: \texttt{POST /upload}
  \item Max size: 50\,MB (configurable)
  \item Storage: \texttt{/data/uploads/<uuid>/<filename>}
  \item Integrity: SHA-256 computed on upload, stored in artifact metadata
  \item Download: \texttt{GET /uploads/<artifact\_id>/<filename>}
\end{itemize}

%----------------------------------------------------------------------
\section{CSV Import/Export}
\begin{itemize}
  \item \texttt{GET /projects/\{id\}/targets/export.csv} --- export targets
  \item \texttt{POST /projects/\{id\}/targets/import} --- import CSV
  \item Columns: name, kind, vendor, model, firmware\_version, environment, notes, tags
\end{itemize}

%----------------------------------------------------------------------
\section{CI/CD Pipeline}
\begin{tabular}{ll}
\toprule
\textbf{Workflow} & \textbf{Triggers} \\
\midrule
\texttt{ci.yml}      & push to main/develop, PRs \\
\texttt{release.yml} & tag \texttt{v*} \\
\bottomrule
\end{tabular}

Jobs: API pytest, LaTeX PDF build, Docker build matrix, Compose smoke test,
GHCR image push, GitHub Release with PDF artifacts.

%----------------------------------------------------------------------
\section{Safety Boundary}
See \texttt{docs/SAFETY\_BOUNDARY.md}.  This platform \textbf{prohibits}:
exploit execution, brute forcing, credential harvesting, malware generation,
evasion workflows, persistence mechanisms, and one-click attacks.
This is NOT a MACOBOX clone, exploit framework, or offensive tool.

%----------------------------------------------------------------------
\section{Raspberry Pi 400 CyberDeck Build}
This platform is designed around the SATUNIX/CYBERDECK hardware BOM:
\begin{itemize}
  \item \textbf{Pi 400} --- keyboard-integrated SBC, Kali Linux ARM 64-bit
  \item \textbf{DFRobot 3.5'' TFT} --- SPI display, waveshare35a dtoverlay
  \item \textbf{TP-Link TL-WN722N v2/v3} --- external WiFi, requires aircrack-ng driver patch
  \item \textbf{Adafruit CyberDeck} --- enclosure and mounting
  \item \textbf{ANKO Battery Pack} --- 5V/2A USB-C portable power
\end{itemize}

Display configuration in \texttt{/boot/config.txt}:
\begin{lstlisting}
dtoverlay=waveshare35a,rotate=270,invertx=1,swapxy=1
\end{lstlisting}

Setup scripts:
\begin{lstlisting}[language=bash]
sudo ./scripts/requirements_setup.sh
sudo ./scripts/lcd_install.sh
sudo ./scripts/tplink_patch.sh
\end{lstlisting}

\vfill
\begin{center}
\textcolor{gray}{\small\copyright\ FLLC \the\year.  For internal authorized use only.}
\end{center}
\end{document}
